\apendice{Especificación de diseño}

\section{Introducción}

Este anexo traduce los requisitos a una arquitectura implementable. Los nombres exactos de clases y tablas podrán adaptarse a las convenciones del código, pero deberán conservarse las responsabilidades, relaciones y restricciones descritas.

\section{Diseño de datos}

\subsection{MySQL como fuente transaccional}

\begin{longtable}{p{0.29\textwidth}p{0.61\textwidth}}
\caption{Entidades transaccionales principales.}\label{tabla:entidades-mysql}\\
\toprule
\textbf{Entidad} & \textbf{Contenido y restricciones}\\
\midrule
\endfirsthead
\toprule
\textbf{Entidad} & \textbf{Contenido y restricciones}\\
\midrule
\endhead
\emph{users} & Identidad, correo normalizado y verificado, nombre, rol, estado y marcas temporales. El correo será único.\\
\emph{software\_apps} & Nombre, slug único, descripción, página oficial, icono, estado, creador y marcas temporales.\\
\emph{tags} y \emph{software\_app\_tags} & Etiquetas únicas y relación muchos a muchos con aplicaciones.\\
\emph{download\_sources} & Aplicación, sistema operativo, arquitectura, región, país, URL de partida, tipo y configuración de resolver, dominios permitidos y estado.\\
\emph{resolved\_sources} & Resultado temporal de resolución: URL protegida, dominio final, nombre, tipo, tamaño, fecha de comprobación y expiración.\\
\emph{bundles} & Propietario opcional para oficiales, nombre, slug, descripción, visibilidad, marca oficial, estado, contador de estrellas y marcas temporales.\\
\emph{bundle\_items} & Relación única entre bundle y aplicación, con preferencias opcionales de sistema y arquitectura.\\
\emph{bundle\_stars} & Relación única usuario-bundle que representa una estrella activa.\\
\emph{download\_jobs} & Usuario o sesión anónima, estado, contexto, tamaños, progreso, referencia temporal al ZIP, errores y expiración.\\
\emph{download\_job\_items} & Aplicación, fuente resuelta, nombre, tamaños, estado, error, reintentos y progreso individual.\\
\emph{software\_requests} & Solicitante, nombre, URL oficial propuesta, estado, respuesta administrativa y fechas.\\
\emph{resolver\_logs} & Fuente, fase, estado, mensaje, metadatos seguros y fecha.\\
\emph{admin\_action\_logs} & Administrador, acción, tipo e identificador del objetivo, resumen del cambio y fecha.\\
\emph{api\_keys} & Usuario, prefijo visible, hash del secreto, estado, última utilización y fecha de revocación.\\
\bottomrule
\end{longtable}

Los datos de configuración que puedan contener secretos no se almacenarán directamente en las tablas. Se utilizarán referencias a secretos externos. Las URL temporales podrán cifrarse o limitarse al almacenamiento del trabajo y nunca se enviarán a registros generales.

\subsection{Estados}

Las aplicaciones utilizarán los estados \emph{active}, \emph{pending review}, \emph{broken}, \emph{disabled}, \emph{deprecated} y \emph{requires manual review}. Solo \emph{active} aparecerá como descargable.

Los trabajos utilizarán \emph{queued}, \emph{resolving}, \emph{checking sizes}, \emph{downloading}, \emph{compressing}, \emph{ready}, \emph{failed}, \emph{cancelled} y \emph{expired}. Las transiciones válidas se comprobarán en un servicio de dominio; no se admitirán saltos arbitrarios desde controladores o workers.

\subsection{Restricciones e índices}

\begin{itemize}
	\item Índices únicos para slugs, correos, etiqueta normalizada, estrella usuario-bundle y elemento bundle-aplicación.
	\item Índices compuestos para fuentes por aplicación, sistema, arquitectura, región y estado.
	\item Índices para trabajos por propietario, estado, creación y expiración.
	\item Índices para bundles públicos por estrellas, fecha y slug.
	\item Integridad referencial con eliminación restringida cuando una aplicación participe en trabajos o registros históricos.
	\item Borrado lógico para aplicaciones y bundles cuando sea necesario conservar auditoría.
\end{itemize}

\subsection{PostgreSQL y pgvector}

PostgreSQL contendrá una tabla derivada similar a:

\begin{verbatim}
software_embeddings
- software_app_id
- content_version
- model_id
- embedding
- generated_at
\end{verbatim}

La clave lógica combinará aplicación, versión del contenido y modelo. Cuando cambien los metadatos relevantes se publicará un trabajo de regeneración. Los resultados semánticos devolverán identificadores; la API consultará MySQL para obtener los datos vigentes y aplicar permisos y filtros.

\section{Diseño arquitectónico}

\subsection{Capas del backend}

\begin{description}
	\item[API:] controladores REST, autenticación, validación de DTO y documentación OpenAPI.
	\item[Aplicación:] casos de uso, transacciones, autorización por propietario y coordinación.
	\item[Dominio:] entidades, reglas, estados, políticas de bundles y objetos de valor.
	\item[Infraestructura:] JPA, RabbitMQ, MinIO, Playwright, correo, embeddings y LLM.
\end{description}

Las dependencias apuntarán hacia las capas interiores. Los adaptadores externos implementarán interfaces definidas por la aplicación o el dominio, especialmente para resolución, almacenamiento temporal, colas y proveedores de IA.

\subsection{Servicios desplegables}

\begin{table}[H]
\centering
\small
\begin{tabularx}{\textwidth}{p{0.23\textwidth}p{0.30\textwidth}X}
\toprule
\textbf{Componente} & \textbf{Responsabilidad} & \textbf{Escalado}\\
\midrule
Frontend & Interfaz estática y consumo de API & Réplicas o CDN\\
API Spring & Usuarios, catálogo, bundles, trabajos y administración & Por tráfico HTTP\\
Worker de resolución & Playwright y validación de fuentes & Por cola de resolución\\
Worker de descarga & Transferencia, progreso y reintentos & Por red y cola\\
Worker de ZIP & Compresión, manifiesto y publicación temporal & Por CPU, disco y cola\\
Servicio semántico & Embeddings, pgvector y coordinación del LLM & Por latencia y cuota\\
Scheduler & Comprobaciones y limpieza periódica & Una instancia con elección de líder\\
\bottomrule
\end{tabularx}
\caption{Componentes lógicos y criterio de escalado.}
\label{tabla:componentes-arquitectura}
\end{table}

La primera implementación podrá combinar los workers en un único proceso. Los mensajes y puertos se diseñarán de forma que su separación posterior no cambie la API pública.

\subsection{TODO: Contratos HTTP principales}

\small
\begin{longtable}{p{0.23\textwidth}p{0.34\textwidth}p{0.33\textwidth}}
\caption{Operaciones principales de la API.}\label{tabla:api-principal}\\
\toprule
\textbf{Método} & \textbf{Ruta conceptual} & \textbf{Finalidad}\\
\midrule
\endfirsthead
\toprule
\textbf{Método} & \textbf{Ruta conceptual} & \textbf{Finalidad}\\
\midrule
\endhead
GET & \emph{/api/apps} & Buscar y filtrar el catálogo.\\
GET & \emph{/api/apps/\{slug\}} & Consultar detalle y fuentes visibles.\\
POST & \emph{/api/download-jobs} & Validar selección y crear trabajo.\\
GET & \emph{/api/download-jobs/\{id\}} & Consultar estado y progreso.\\
DELETE & \emph{/api/download-jobs/\{id\}} & Cancelar un trabajo.\\
GET & \emph{/api/download-jobs/\{id\}/file} & Descargar un ZIP listo.\\
GET/POST & \emph{/api/bundles} & Buscar o crear bundles.\\
GET/PATCH/DELETE & \emph{/api/bundles/\{slug\}} & Consultar o mantener un bundle propio.\\
POST & \emph{/api/bundles/\{id\}/star} & Asignar estrella.\\
DELETE & \emph{/api/bundles/\{id\}/star} & Retirar estrella.\\
POST & \emph{/api/software-requests} & Solicitar una aplicación.\\
POST & \emph{/api/semantic-search} & Buscar por intención.\\
POST & \emph{/api/bundle-suggestions} & Crear un borrador de bundle.\\
POST & \emph{/api/admin/sources/\{id\}/test} & Probar una fuente y su resolver.\\
\bottomrule
\end{longtable}
\normalsize

Los cuerpos y respuestas se definirán en OpenAPI. Los errores compartirán un esquema con código estable, mensaje público, identificador de incidencia y detalles de validación permitidos.

\section{Diseño procedimental}

\subsection{Creación de un trabajo}

\begin{enumerate}
	\item Autenticar o identificar la sesión anónima.
	\item Validar lista, sistema, arquitectura, región y formato.
	\item Eliminar duplicados y comprobar que las aplicaciones están activas.
	\item Aplicar cuotas y límites conocidos.
	\item Crear \ruta{download\_jobs} y sus elementos en una transacción.
	\item Publicar el identificador en RabbitMQ mediante patrón de salida transaccional o mecanismo equivalente.
	\item Responder con código de aceptación e identificador del trabajo.
\end{enumerate}

\subsection{Resolución segura}

\begin{enumerate}
	\item Seleccionar la fuente compatible más específica.
	\item Comprobar si existe una resolución no expirada.
	\item Ejecutar resolver dedicado o Playwright.
	\item Validar HTTPS, hostname, DNS, IP y puerto.
	\item Seguir un número limitado de redirecciones repitiendo la validación.
	\item Comprobar dominio final, tipo de contenido, nombre y tamaño.
	\item Persistir resultado y metadatos sin exponer secretos.
	\item Marcar para revisión manual cuando el fallo sea permanente.
\end{enumerate}

\subsection{Descarga y reintentos}

Cada elemento se descargará con un máximo configurable de intentos. Los fallos temporales utilizarán espera incremental. Si la fuente acepta rangos, una implementación posterior podrá reanudar; en la primera versión se reiniciará el archivo para simplificar la validación.

Durante la transferencia se controlará el número real de bytes. Al terminar se comparará con el tamaño anunciado cuando exista, se verificará el tipo y se rechazará cualquier archivo que exceda límites o cambie a un dominio no autorizado.

\subsection{Compresión y finalización}

\begin{enumerate}
	\item Reunir elementos válidos y errores finales.
	\item Sanitizar y resolver nombres duplicados.
	\item Generar \ruta{manifest.json} y \ruta{README.txt}.
	\item Construir el ZIP en el espacio temporal.
	\item Publicarlo en MinIO con expiración.
	\item Marcar el trabajo como listo o fallido.
	\item Eliminar instaladores intermedios.
	\item Enviar aviso si fue solicitado.
\end{enumerate}

\subsection{Cancelación y limpieza}

La cancelación fijará una marca persistente y publicará una señal para los workers. Cada fase comprobará esa marca entre operaciones y cerrará recursos. El proceso eliminará objetos parciales y finalizará como \emph{cancelled}. El scheduler buscará trabajos expirados o bloqueados y ejecutará una limpieza idempotente.

\subsection{Búsqueda semántica}

\begin{enumerate}
	\item Validar longitud y filtros de la consulta.
	\item Generar embedding mediante el proveedor configurado.
	\item Recuperar los vecinos más próximos en pgvector.
	\item Consultar las aplicaciones actuales en MySQL.
	\item Aplicar estado, compatibilidad, seguridad y permisos.
	\item Devolver resultados vectoriales.
	\item Si se solicita explicación, enviar solo esos resultados al LLM.
	\item Validar el esquema y descartar referencias no presentes.
\end{enumerate}

\subsection{Actualización del contador de estrellas}

La estrella se insertará con una restricción única. El contador del bundle se actualizará en la misma transacción o mediante un evento idempotente. Una tarea administrativa podrá recalcularlo desde \ruta{bundle\_stars} para reparar discrepancias.

\section{Diseño de seguridad}

Los secretos se proporcionarán mediante el entorno o un gestor de secretos. Los contenedores se ejecutarán sin privilegios, con sistemas de archivos de solo lectura cuando sea posible y límites de CPU, memoria y almacenamiento efímero.

Los logs utilizarán campos estructurados y excluirán contraseñas, claves, cabeceras de autorización, URL firmadas y contenido de instaladores. La auditoría registrará cambios de seguridad sin copiar secretos.

La salida del LLM se considerará entrada no confiable: no se renderizará como HTML sin sanitización, no se ejecutará como código y no se convertirá directamente en comandos o cambios de configuración.

\section{Diseño de despliegue}

El entorno local utilizará Docker Compose. El despliegue escalable utilizará Docker con:

\begin{itemize}
	\item \emph{Deployments} separados para API y workers;
	\item \emph{Services} internos para componentes sin exposición pública;
	\item \emph{Ingress} únicamente para frontend y API;
	\item comprobaciones de vida y disponibilidad;
	\item límites y solicitudes de recursos;
	\item volúmenes o almacenamiento de objetos para temporales controlados;
	\item secretos fuera de las imágenes;
	\item políticas de red que restrinjan el acceso entre componentes.
\end{itemize}

El escalado automático de workers se basará preferentemente en la longitud de las colas y en métricas de recursos. No se escalará Playwright sin límites, porque cada navegador consume memoria y conexiones externas.
